\begin{recipe}{Couscous met feta, groenten en rozijntjes}
\kicker{Hoofdgerecht · vegetarisch · Mediterraans}
\index[register]{Couscous met feta, groenten en rozijntjes}
\index[register]{Couscous}
\index[register]{Feta}
\index[register]{Spinazie}

\meta{20-25 minuten \dvd Mediterraans}

\lettrine{E}{en} snelle maaltijdsalade met warme paprika en spinazie, romige
feta en een zoet accent van rozijnen en pitten. Fijn doordeweeks gerecht, en
handig als de kinderen nog moeten wennen aan de smaak van feta: die kun je
apart bij de couscous zetten in plaats van er meteen doorheen te roeren.

\blockrule{Ingrediënten}
\begin{ingredients}
  \ing{420 g}{volkoren couscous}\index[register]{Couscous}
  \ing{1}{ui}\index[register]{Ui}
  \ing{2½}{gele paprika}\index[register]{Paprika}
  \ing{1}{citroen}
  \ing{625 g}{spinazie}\index[register]{Spinazie}
  \ing{100 g}{rozijnen-pittenmix}\index[register]{Rozijnen}
  \ing{150 g}{feta}\index[register]{Feta}
  \ing{850 ml}{groentebouillon}
  \ingb{olijfolie}
  \ing{5 tl}{oregano}
  \ingb{versgemalen peper en zout}
\end{ingredients}

\blockrule{Bereiding}
\begin{steps}
  \item Breng de groentebouillon aan de kook. Meng de couscous met de
        bouillon in een kom en laat, afgedekt, 10 minuten wellen. Roer daarna
        los met een vork.
  \item Snipper ondertussen de ui. Snijd de paprika grof. Rasp de schil
        van de citroen en pers het sap uit.
  \item Verhit olijfolie in een wok of hapjespan met deksel en fruit de
        ui en oregano 1 minuut op laag vuur. Voeg de paprika en het
        citroenrasp toe, draai het vuur middelhoog en bak 4 minuten. Voeg
        daarna de spinazie toe, eventueel in delen, en laat slinken.
  \item Voeg de couscous en het citroensap toe aan de wok of hapjespan.
        Breng op smaak met peper en zout en roerbak nog 1 minuut. Schep de
        helft van de rozijnen-pittenmix erdoor.
  \item Verdeel de couscous met groenten over de borden. Bestrooi met de
        overige rozijnen-pittenmix en verkruimel de feta erover. Besprenkel
        naar smaak met extra citroensap en olijfolie.
\end{steps}
\end{recipe}
